\subsection{Upper-Triangular Matrices}

\begin{exercise}{5c1}
    Prove or give a counterexample: If $T \in \lin{V}$ and $T^2$ has an upper-triangular matrix with respect to some basis of $V$, then $T$ has an upper-triangular matrix with respect to some basis of $V$.

    \remark{The matrix in this exercise represents the $90^\circ$ counterclockwise rotation of $\rbb^2$.}
\end{exercise}

\begin{sol}
    We claim that the statement is false. 
    Consider the linear operator $T \in \lin{\rbb^2}$ with matrix representation
    \[
        A = 
        \begin{pmatrix}
            0 & 1 \\
            -1 & 0
        \end{pmatrix}
    \]
    with respect to the standard basis. 
    Then $T^2$ is given by
    \[
        A^2 = 
        \begin{pmatrix}
            0 & 1 \\
            -1 & 0
        \end{pmatrix}^2 = 
        \begin{pmatrix}
            -1 & 0 \\
            0 & -1
        \end{pmatrix},
    \]
    which is an upper-triangular matrix with respect to the standard basis of $\rbb^2$. 
    However, $T$ does not have an upper-triangular matrix with respect to any basis of $\rbb^2$, since it has no real eigenvalues and thus cannot be represented in upper-triangular form over $\rbb$.
\end{sol}

\begin{exercise}{5c2}
    Suppose $A$ and $B$ are upper-triangular matrices of the same size, with $\alpha_1, \ldots, \alpha_n$ on the diagonal of $A$ and $\beta_1, \ldots, \beta_n$ on the diagonal of $B$.
    \begin{subexercises}
        \item Show that $A + B$ is an upper-triangular matrix with $\alpha_1 + \beta_1, \ldots, \alpha_n + \beta_n$ on the diagonal.
        \item Show that $AB$ is an upper-triangular matrix with $\alpha_1\beta_1, \ldots, \alpha_n\beta_n$ on the diagonal.
    \end{subexercises}
    \note{The results in this exercise are used in the proof of 5.81.}
\end{exercise}

\begin{solenum}
    \item Let $C = A + B$.
    Since $A$ and $B$ are upper-triangular, for any entry $c_{ij}$ of $C$ where $i > j$, we have
    \[
        c_{ij} = a_{ij} + b_{ij} = 0 + 0 = 0,
    \]
    which shows that $C$ is upper-triangular.
    The diagonal entries of $C$ are given by
    \[
        c_{ii} = a_{ii} + b_{ii} = \alpha_i + \beta_i,
    \]
    for $i = 1, \ldots, n$.
    \item Let $D = AB$. 
    Since $A$ and $B$ are upper-triangular, for any entry $d_{ij}$ of $D$ where $i > j$, we have the following sum for the $(i,j)$-th entry of $D$:
    \[
        d_{ij} = \sum_{k=1}^n a_{ik}b_{kj}.
    \]
    Consider $a_{ik}$ when $k < i$. Since $A$ is upper-triangular, $a_{ik} = 0$.
    Now consider $b_{kj}$ when $k > j$. Since $B$ is upper-triangular, $b_{kj} = 0$.
    Therefore, for $i > j$, all terms in the sum are zero, and we have $d_{ij} = 0$ for $i > j$, showing that $D$ is upper-triangular.

    To see that the diagonal entries of $D$ are $\alpha_1\beta_1, \ldots, \alpha_n\beta_n$, we compute the diagonal entries:
    \[
        d_{ii} = \sum_{k=1}^n a_{ik}b_{ki},
    \]
    for the $i$-th diagonal entry. Since $A$ and $B$ are upper-triangular, then $k < i$ implies $a_{ik} = 0$ and $k > i$ implies $b_{ki} = 0$. Therefore, the only nonzero term in the sum occurs when $k = i$, giving
    \[
        d_{ii} = a_{ii}b_{ii} = \alpha_i\beta_i,
    \]
    for $i = 1, \ldots, n$.
\end{solenum}

\begin{exercise}{5c3}
    Suppose $T \in \lin{V}$ is invertible and $v_1, \ldots, v_n$ is a basis of $V$ with respect to which the matrix of $T$ is upper triangular, with $\lambda_1, \ldots, \lambda_n$ on the diagonal.
    Show that the matrix of $T\inv$ is also upper triangular with respect to the basis $v_1, \ldots, v_n$, with
    \[
        \frac{1}{\lambda_1}, \ldots, \frac{1}{\lambda_n}
    \]
    on the diagonal.

    \apnote{5c3}{ for a direct computation of the entries of $T\inv$.}
\end{exercise}

\begin{sol}
    Let $V_k=\spn(v_1, \dots, v_k)$ for each $k \in \{1, \dots, n\}$.
    Because $\mcal(T,(v_1,\dots,v_n))$ is upper-triangular, $T(V_k) \subseteq V_k$ for every $k$; thus each $V_k$ is invariant under $T$.
    Since $T$ is invertible, its restriction to the (finite-dimensional) $V_k$ is injective, hence surjective onto $V_k$.
    Therefore, $T(V_k)=V_k$. Applying $T\inv$ yields $V_k=T\inv (V_k)$, meaning each $V_k$ is also invariant under $T\inv$. So, $\mcal(T\inv, (v_1, \dots, v_n))$ is upper-triangular. 

    Let $\mu_1$, $\dots$, $\mu_n$ be the diagonal entries of $\mcal(T\inv, (v_1, \dots, v_n))$.
    Since $TT\inv=I$ is upper-triangular, we have from \exref{5c2} (b) that $\lambda_k \mu_k=1$ for each $k$, giving $\mu_k=\frac{1}{\lambda_k}$.
\end{sol}

\begin{exercise}{5c4}
    Give an example of an operator whose matrix with respect to some basis contains only $0$'s on the diagonal, but the operator is invertible.

    \note{This exercise and the exercise below show that 5.41 fails without the hypothesis that an upper-triangular matrix is under consideration.}
\end{exercise}

\begin{sol}
    Consider the linear operator $T \in \lin{\rbb^2}$ with matrix representation
    \[
        A = 
        \begin{pmatrix}
            0 & 1 \\
            -1 & 0
        \end{pmatrix}
    \]
    with respect to the standard basis.
    $A$ is invertible with inverse
    \[
        A\inv = 
        \begin{pmatrix}
            0 & -1 \\
            1 & 0
        \end{pmatrix}. \qedhere
    \]
\end{sol}

\begin{exercise}{5c5}
    Give an example of an operator whose matrix with respect to some basis contains only nonzero numbers on the diagonal, but the operator is not invertible.
\end{exercise}

\begin{sol}
    Consider the linear operator $T \in \lin{\rbb^2}$ with matrix representation
    \[
        A = 
        \begin{pmatrix}
            1 & 1 \\
            -1 & -1
        \end{pmatrix}
    \]
    with respect to the standard basis.
    Since $T(1, -1) = (0, 0)$, $T$ is not invertible.
\end{sol}

\begin{exercise}{5c6}
    Suppose $\fbb = \cbb$, $V$ is finite-dimensional, and $T \in \lin{V}$.
    Prove that if $k \in \{1, \ldots, \dim V\}$, then $V$ has a $k$-dimensional subspace invariant under $T$.
\end{exercise}

\begin{sol}
    Recall that every operator on a finite-dimensional complex vector space has an upper-triangular matrix with respect to some basis.
    Let $v_1$, $\dots$, $v_n$ be such a basis.
    Then, for each $k \in \{1, \dots, n\}$,
    \[
        Tv_j \in \spn(v_1, \dots, v_j) \subseteq \spn(v_1, \dots, v_k)
    \]
    for all $j \leq k$.
    Hence, $\spn(v_1, \dots, v_k)$ is a $k$-dimensional subspace that is invariant under $T$.
\end{sol}

\begin{exercise}{5c7}
    Suppose $V$ is finite-dimensional, $T \in \lin{V}$, and $v \in V$.
    \begin{subexercises}
        \item Prove that there exists a unique monic polynomial $p_v$ of smallest degree such that $p_v(T)v = 0$.
        \item Prove that the minimal polynomial of $T$ is a polynomial multiple of $p_v$.
    \end{subexercises}
\end{exercise}

\begin{solenum}
    \item If $v = 0$, then $p(T)v = 0$ for any polynomial $p$, so we can take $p_v(x) = 1$.
    Otherwise, suppose $v \neq 0$.
    Consider the list $v, Tv, T^2v, \ldots, T^{\dim V}v$.
    Since $V$ is finite-dimensional, this list is linearly dependent.
    Then there exist scalars $a_0, \ldots, a_m$ not all zero such that
    \[
        a_0v + a_1Tv + \cdots + a_mT^mv = 0.
    \]
    Here, $m$ denotes the smallest integer such that there exist scalars $a_0, \ldots, a_m$ not all zero satisfying the above equation.
    Since $m$ is minimal, we have $a_m \neq 0$.
    Let $p_v(x)$ be the monic polynomial after dividing by $a_m$:
    \[
        p_v(x) = x^m + \frac{a_{m-1}}{a_m}x^{m-1} + \cdots + \frac{a_0}{a_m}.
    \]
    Then $p_v(T)v = 0$.
    Suppose $q$ is another monic polynomial of smallest degree such that $q(T)v = 0$.
    Then $p_v - q$ is a polynomial of smaller degree than $p_v$ such that $(p_v - q)(T)v = 0$, which contradicts the minimality of $p_v$.
    \item Let $m_T$ be the minimal polynomial of $T$.
    Then $m_T(T) = 0$, so $m_T(T)v = 0$.
    By the division algorithm for polynomials, there exist unique polynomials $q'$ and $r$ such that
    \[
        m_T = p_vq' + r,
    \]
    where $r = 0$ or $\deg r < \deg p_v$.
    Then
    \[
        0 = m_T(T)v = (p_vq' + r)(T)v = p_v(T)q'(T)v + r(T)v = q'(T)p_v(T)v + r(T)v = r(T)v,
    \]
    since $p_v(T)v = 0$.
    If $r \neq 0$, then $r$ is a polynomial of smaller degree than $p_v$ such that $r(T)v = 0$, which contradicts the minimality of $p_v$.
    Therefore, $r = 0$, and we have $m_T = p_vq'$, which shows that $m_T$ is a polynomial multiple of $p_v$.
\end{solenum}

\begin{exercise}{5c8}
    Suppose $V$ is finite-dimensional, $T \in \lin{V}$, and there exists a nonzero vector $v \in V$ such that $T^2v + 2Tv = -2v$.
    \begin{subexercises}
        \item Prove that if $\fbb = \rbb$, then there does not exist a basis of $V$ with respect to which $T$ has an upper-triangular matrix.
        \item Prove that if $\fbb = \cbb$ and $A$ is an upper-triangular matrix that equals the matrix of $T$ with respect to some basis of $V$, then $-1 + i$ or $-1 - i$ appears on the diagonal of $A$.
    \end{subexercises}
\end{exercise}

\begin{solenum}
    \item Suppose there exists a basis of $V$ with respect to which $T$ has an upper-triangular matrix.
    Then the diagonal entries of this matrix are the eigenvalues of $T$.
    Let $p(z) = z^2 + 2z + 2$.
    Then $p(T)$ also has an upper-triangular matrix with respect to the same basis, and the diagonal entries of this matrix are $p(\lambda)$, where $\lambda$ is an eigenvalue of $T$.
    Since $p(T)v = 0$ for some nonzero vector $v$, $p(T)$ is not injective.
    Then $p(T)$ is not invertible, so $p(\lambda) = 0$ for some eigenvalue $\lambda \in \rbb$ of $T$.
    However, $p(z)$ has no real roots, contradicting that $p(\lambda) = 0$ for some real eigenvalue $\lambda$.
    \item Utilizing the same deductions as in part (a), we conclude that there exists an eigenvalue $\lambda$ of $T$ such that $p(\lambda) = 0$.
    The roots of $p(z)$ are $-1 + i$ and $-1 - i$, so either $-1 + i$ or $-1 - i$ must be an eigenvalue of $T$.
    Since the diagonal entries of an upper-triangular matrix are the eigenvalues of the corresponding operator, either $-1 + i$ or $-1 - i$ must appear on the diagonal of $A$.
\end{solenum}

\begin{exercise}{5c9}
    Suppose $B$ is a square matrix with complex entries.
    Prove that there exists an invertible square matrix $A$ with complex entries such that $A\inv BA$ is an upper-triangular matrix.
\end{exercise}

\begin{sol}
    Since $B$ represents a linear operator on a finite-dimensional complex vector space, let $T$ be the linear operator corresponding to $B$, i.e., $B = \mcal(T, (e_1, \ldots, e_n))$, where $(e_1, \ldots, e_n)$ is the standard basis of $\cbb^n$.
    Since $T$ has an upper-triangular matrix with respect to some basis of the vector space, let $A$ be the change of basis matrix from the new basis to the standard basis.
    Then $A\inv BA$ is the matrix representation of $T$ with respect to this new basis, which is upper-triangular.
\end{sol}

\begin{exercise}{5c10}
    Suppose $T \in \lin{V}$ and $v_1, \ldots, v_n$ is a basis of $V$.
    Show that the following are equivalent.
    \begin{subexercises}
        \item The matrix of $T$ with respect to $v_1, \ldots, v_n$ is lower triangular.
        \item $\spn(v_k, \ldots, v_n)$ is invariant under $T$ for each $k = 1, \ldots, n$.
        \item $Tv_k \in \spn(v_k, \ldots, v_n)$ for each $k = 1, \ldots, n$.
    \end{subexercises}
    \note{A square matrix is called \textbf{lower triangular} if all entries above the diagonal are $0$.}
\end{exercise}

\begin{solitem}
    \item Suppose the matrix of $T$ with respect to $v_1, \ldots, v_n$ is lower triangular.
    Fix some $k \in \set{1, \ldots, n}$.
    For $j \in \set{1, \ldots, n}$, then
    \[
        Tv_j \in \spn(v_j, \ldots, v_n)
    \]
    because the matrix is lower triangular.
    If $j \geq k$, then $\spn(v_j, \ldots, v_n) \subseteq \spn(v_k, \ldots, v_n)$.
    Therefore, $Tv_j \in \spn(v_k, \ldots, v_n)$ for all $j \geq k$, which shows that $\spn(v_k, \ldots, v_n)$ is invariant under $T$.
    \item Suppose $\spn(v_k, \ldots, v_n)$ is invariant under $T$ for each $k = 1, \ldots, n$.
    Then for each $k$, we have $Tv_k \in \spn(v_k, \ldots, v_n)$, which shows that $Tv_k$ can be expressed as a linear combination of $v_k, \ldots, v_n$.
    \item Suppose $Tv_k \in \spn(v_k, \ldots, v_n)$ for each $k = 1, \ldots, n$.
    Then
    \[
        Tv_k = \sum_{j = k}^n A_{j, k}v_j
    \]
    for some scalars $A_{j, k}$.
    Then $A_{i, k} = 0$ for all $i < k$, which shows that the matrix of $T$ with respect to $v_1, \ldots, v_n$ is lower triangular.
\end{solitem}

\begin{exercise}{5c11}
    Suppose $\fbb = \cbb$ and $V$ is finite-dimensional.
    Prove that if $T \in \lin{V}$, then there exists a basis of $V$ with respect to which $T$ has a lower-triangular matrix.
\end{exercise}

\begin{sol}
    Since $T$ has an upper-triangular matrix with respect to some basis of $V$, let $v_1, \ldots, v_n$ be such a basis.
    Then $\spn(v_1, \ldots, v_k)$ is invariant under $T$ for each $k = 1, \ldots, n$.
    Consider the reverse basis $w_1 = v_n, w_2 = v_{n-1}, \ldots, w_n = v_1$.
    Then $\spn(w_k, \ldots, w_n) = \spn(v_1, \ldots, v_{n-k+1})$ is invariant under $T$ for each $k = 1, \ldots, n$.
    By \exref{5c10}, the matrix of $T$ with respect to the basis $w_1, \ldots, w_n$ is lower triangular.
\end{sol}

\begin{exercise}{5c12}
    Suppose $V$ is finite-dimensional, $T \in \lin{V}$ has an upper-triangular matrix with respect to some basis of $V$, and $U$ is a subspace of $V$ that is invariant under $T$.
    \begin{subexercises}
        \item Prove that $T|_U$ has an upper-triangular matrix with respect to some basis of $U$.
        \item Prove that the quotient operator $T/U$ has an upper-triangular matrix with respect to some basis of $V/U$.
    \end{subexercises}
    \note{The quotient operator $T/U$ was defined in \exref{5a38}.}
\end{exercise}

\begin{solenum}
    \item Since $T$ has an upper-triangular matrix with respect to some basis of $V$, then its minimal polynomial is a product of linear factors.
    Since $U$ is invariant under $T$, the minimal polynomial of $T|_U$ divides the minimal polynomial of $T$, and thus is also a product of linear factors.
    Therefore, $T|_U$ has an upper-triangular matrix with respect to some basis of $U$.
    \item Since the minimal polynomial of $T$ is a polynomial multiple of the minimal polynomial of $T/U$, the same argument as in part (a) shows that $T/U$ has an upper-triangular matrix with respect to some basis of $V/U$.
\end{solenum}

\begin{exercise}{5c13}
    Suppose $V$ is finite-dimensional and $T \in \lin{V}$.
    Suppose there exists a subspace $U$ of $V$ that is invariant under $T$ such that $T|_U$ has an upper-triangular matrix with respect to some basis of $U$ and also $T/U$ has an upper-triangular matrix with respect to some basis of $V/U$.
    Prove that $T$ has an upper-triangular matrix with respect to some basis of $V$.
\end{exercise}

\begin{sol}
    Let $u_1, \ldots, u_m$ be a basis of $U$ with respect to which $T|_U$ has an upper-triangular matrix.
    Let $v_1 + U, \ldots, v_n + U$ be a basis of $V/U$ with respect to which $T/U$ has an upper-triangular matrix.
    Then the list $u_1, \ldots, u_m, v_1, \ldots, v_n$ is a basis of $V$.
    Since $T|_U$ has an upper-triangular matrix with respect to the basis $u_1, \ldots, u_m$, we have
    \[
        Tu_j = \sum_{k=1}^j A_{k,j}u_k
    \]
    for some scalars $A_{k,j}$ and for each $j = 1, \ldots, m$.
    Since $T/U$ has an upper-triangular matrix with respect to the basis $v_1 + U, \ldots, v_n + U$, we have
    \[
        T(v_j + U) = \sum_{k=1}^j B_{k,j}(v_k + U)
    \]
    for some scalars $B_{k,j}$ and for each $j = 1, \ldots, n$.
    This implies that
    \[
        Tv_j - \sum_{k=1}^j B_{k,j}v_k \in U,
    \]
    so there exist scalars $C_{l,j}$ such that
    \[
        Tv_j - \sum_{k=1}^j B_{k,j}v_k = \sum_{l=1}^m C_{l,j}u_l.
    \]
    Therefore,
    \[
        Tv_j = \sum_{l=1}^m C_{l,j}u_l + \sum_{k=1}^j B_{k,j}v_k,
    \]
    which shows that the matrix of $T$ with respect to the basis $u_1, \ldots, u_m, v_1, \ldots, v_n$ is upper-triangular.
\end{sol}

\begin{exercise}{5c14}
    Suppose $V$ is finite-dimensional and $T \in \lin{V}$.
    Prove that $T$ has an upper-triangular matrix with respect to some basis of $V$ if and only if the dual operator $T'$ has an upper-triangular matrix with respect to some basis of the dual space $V'$.
\end{exercise}

\begin{sol}
    By \exref{5b28}, the minimal polynomials of $T$ and $T'$ are the same.
    Therefore, the following are equivalent:
    \begin{itemize}
        \item $T$ has an upper-triangular matrix with respect to some basis of $V$.
        \item The minimal polynomial of $T$ is a product of linear factors.
        \item The minimal polynomial of $T'$ is a product of linear factors.
        \item $T'$ has an upper-triangular matrix with respect to some basis of $V'$. \qedhere
    \end{itemize}
\end{sol}